
\begin{proposition}[Every irreducible class VII normalization contributes zero]\label{cdp:st:C.2.MH.4}
Under the full CDP torus-fibration hypotheses, an
irreducible fibre support whose normalization resolution
is a surface of class VII has
\[
                              e(S)=0.
 \tag{C.MH.16}\label{cdp:eq:C.MH.16}
\]
Consequently every such actual fibre satisfies the
equality data of \ref{cdp:st:C.2.VII.4}.
\end{proposition}
\begin{proof}
Use the prime-order rational augmentation local system
$V$ of \ref{cdp:st:C.2.MH.2}, of rank $d=\ell-1$. Its construction,
the torus acyclicity (C.MH.8), and the support duality
(C.MH.13) do not require a projective normalization.
They give $h^0(S,V)=h^1(S,V)=0$ and
$h^2(S,V)=h^4(S,V)$. Conductor descent still gives the
surjection $H^3(S,V)\to H^3(T,V)$ and the isomorphism
$H^4(S,V)\simeq H^4(T,V)$. The negative-definite
exceptional Gysin matrix in the proof of (C.MH.12)
is valid for the resolution of any normal complex
surface. It gives
\[
 \begin{gathered}
 h^3(S,V)\ge h^3(Z,V),\qquad h^4(S,V)=h^4(Z,V),\\
 d\,e(S)\le 2h^4(Z,V)-h^3(Z,V).
 \end{gathered}
 \tag{C.MH.17}\label{cdp:eq:C.MH.17}
\]
No mixed Hodge assertion is used in these comparisons.

If the cyclic character restricts nontrivially to $Z$,
all its nontrivial powers remain nontrivial because
$\ell$ is prime. Poincar\'e duality on the compact
oriented surface gives $h^4(Z,V)=0$, and \eqref{cdp:eq:C.MH.17}
gives $e(S)\le0$. If its restriction is trivial,
$h^4(Z,V)=d$ and $h^3(Z,V)=d$: a class VII surface
has $b_1=b_3=1$, and point blowups preserve $b_3$.
Thus $e(S)\le1$ in this second case. \ref{cdp:st:C.2.UN.1} proves
that the Euler number of every actual torus-degeneration
fibre is even, so the second case also gives
$e(S)\le0$. The nonnegative lower bound of \ref{cdp:st:C.2.VII.3}
now gives \eqref{cdp:eq:C.MH.16} in both cases. \ref{cdp:st:C.2.VII.4} supplies
its exact normalization-and-conductor equality data.
\end{proof}


